\section{Real-World Robot Validation}
\label{sec:hardware}

Simulation and mathematical analysis provide controlled access to internal physical quantities, but the final objective is to design robots that can be built and operated. This section is reserved for the physical-robot demonstrations and measurements that have already been organized for the project. They will be integrated as a separate validation layer rather than being mixed with the eleven simulation FULL cases.

\subsection{Robot platform and experimental setup}

\todo{Insert the physical robot platform, morphology, tendon routing, materials, actuation/sensing arrangement, calibration procedure, and the correspondence between the simulated and manufactured structures.}

\subsection{Demonstration tasks and protocol}

\todo{Select the physical tasks to be shown. For each task, specify the initial condition, target or loading condition, execution protocol, number of trials, and the metrics used for evaluation.}

\subsection{Physical results}

\todo{Insert the already archived real-world measurements, photographs, and demonstration frames. Report measured values and trial counts from the source records only; do not infer hardware performance from simulation results.}

\begin{figure*}[t]
\centering
\fbox{\parbox[c][3.0cm][c]{0.94\linewidth}{\centering
\todo{Insert selected physical-robot photographs and synchronized demonstration frames. Keep the task identity and experimental condition visible in each panel.}}}
\caption{Physical-robot demonstrations. The final figure will use the existing real-world demonstration records selected for the paper.}
\label{fig:real-world}
\end{figure*}

The physical-robot section is intentionally kept separate from the simulation task portfolio. The final manuscript will report the hardware configuration, protocol, and measured outcomes from the corresponding experimental records.
